\section{Neighbourhoods of the cubic lattices}\label{sec:local}

Near a point of $\Kbox_{\mathrm{red}}$ that represents $\ltwo\Bz$ or $\ltwo\Fz$ the deficit vanishes to second order, and no bound that
treats the energy and the defect separately can succeed. In this section we prove the main inequality in explicit neighbourhoods of
all such points, by a Taylor expansion with certified coefficients.

Theorem~\ref{thm:local} is a result of its own: the proof of the main theorems uses neither it nor any other part of this section,
because the cover of Section~\ref{sec:cover} treats these neighbourhoods by the anchor rules of Section~\ref{ssec:anchor}.

\subsection{The fourteen two-site descriptions}\label{ssec:fourteen}

Let $\Lambda_0\in\{\Bz,\Fz\}$ and let $\Lambda_0^\sharp=\{k:k\cdot b\in\Z\text{ for all }b\in\Lambda_0\}$. The sublattices of index two
of $\Lambda_0$ are the sets $L_0(k)=\{b\in\Lambda_0:k\cdot b\in\Z\}$, where $k$ runs over the seven nonzero classes of
$\frac12\Lambda_0^\sharp/\Lambda_0^\sharp$. Geometrically, $L_0(k)$ consists of every second lattice plane perpendicular to $k$. The
point group of $\Lambda_0$, the $48$ signed permutation matrices, permutes these classes, with the following orbits.

\begin{table}[h]
\begin{tabular}{lllc}
\toprule
lattice & class of $k$ & planes & number\\
\midrule
$\Bz$ & $(\frac12,\frac12,0)$ & $\{110\}$ & 6\\
$\Bz$ & $(1,0,0)$ & $\{100\}$ & 1\\
$\Fz$ & $2^{2/3}(\frac14,\frac14,\frac14)$ & $\{111\}$ & 4\\
$\Fz$ & $2^{2/3}(\frac12,0,0)$ & $\{100\}$ & 3\\
\bottomrule
\end{tabular}
\caption{The sublattices of index two of the cubic lattices, up to symmetry. In the language of lattice dynamics the classes of $k$ are
the points $N$ and $H$ of the Brillouin zone of the body-centred cubic lattice and the points $L$ and $X$ of that of the face-centred
cubic lattice.}\label{tab:classes}
\end{table}

For instance $L_0=\Z^3$ for the class $(1,0,0)$ of $\Bz$: this is the familiar description of the body-centred cubic lattice as two
interpenetrating simple cubic lattices. A triple $(\Lambda_0,L_0,w_0)$ with $w_0\in\Lambda_0\setminus L_0$ is called a
\emph{two-site description} of $\Lambda_0$. Up to the choice of $w_0$ in its coset, which is immaterial below, there are fourteen
two-site descriptions, in four classes.

\subsection{Charts}\label{ssec:charts}

Fix a two-site description $(\Lambda_0,L_0,w_0)$. For a traceless symmetric matrix $S$ with $I+S$ positive definite and $u\in\R^3$ let
\begin{equation}\label{eq:chart}
  C(S):=\det(I+S)^{-1/3}(I+S),\qquad A(S):=C(S)^{1/2},\qquad Y(S,u):=A(S)\bigl(L_0\cup(L_0+w_0+u)\bigr).
\end{equation}
Thus $S$ parametrizes the positive definite matrices $C$ of determinant one, with inverse $S=3C/\tr C-I$; $A(S)$ is the symmetric
deformation with Gram matrix $C$; and $u$ is the displacement of the second sublattice in the undeformed frame. We use the Frobenius
norm $\|S\|_\Fro$ on the five-dimensional space of traceless symmetric matrices, and for $r,\nu>0$ the ellipsoid
\[
  \mathcal E(r,\nu):=\bigl\{(S,u):\ \|S\|_\Fro^2+|u|^2/\nu^2\le r^2\bigr\}\subset\R^8 .
\]

\begin{lemma}[Chart coordinates of a configuration]\label{lem:membership}
Let $B'$ be a basis matrix of $L_0$ and $f_0:=B'^{-1}w_0$. Let $Y=R\,\Z^3\cup(R\,\Z^3+Rf)$ be a normalized two-site configuration and
$n\in\Z^3$. Put
\[
  C:=B'^{-\T}R^\T RB'^{-1},\qquad S:=\frac{3C}{\tr C}-I,\qquad u:=B'(f-f_0-n).
\]
Then $Y=Q\,Y(S,u)$ for some $Q\in\Or(3)$.
\end{lemma}

\begin{proof}
$C$ is positive definite with $\det C=1$, so $C=C(S)$ and $RB'^{-1}=QA(S)$ with $Q$ orthogonal. Then $R\,\Z^3=QA(S)L_0$ and
$Rf=QA(S)(w_0+u)+Rn$.
\end{proof}

A signed permutation matrix $O$ maps $(\Lambda_0,L_0,w_0)$ to another two-site description, and
$Y_{OL_0,Ow_0}(OSO^\T,Ou)=O\,Y_{L_0,w_0}(S,u)$. The map $(S,u)\mapsto(OSO^\T,Ou)$ preserves each $\mathcal E(r,\nu)$. Hence a statement
about $\Def$ on $\mathcal E(r,\nu)$ that holds for one description holds for all descriptions of its class.

\subsection{The local inequality}

\begin{theorem}[Local inequality]\label{thm:local}
Let $(\Lambda_0,L_0,w_0)$ be any of the fourteen two-site descriptions of $\Bz$ and $\Fz$, and let $(r,\nu)$ be one of the two pairs
listed for its class in Table~\ref{tab:atlas}. Then, with $\varepsilon_0>0$ from the table,
\begin{equation}\label{eq:localineq}
  \Def\bigl(\ltwo Y(S,u)\bigr)\ \ge\ \varepsilon_0\Bigl(\|S\|_\Fro^2+\frac{|u|^2}{\nu^2}\Bigr)\qquad\text{for all }(S,u)\in\mathcal E(r,\nu).
\end{equation}
In particular $\Def>0$ on $\mathcal E(r,\nu)\setminus\{0\}$.
\end{theorem}

Near a two-site description of $\Fz$ the defect is identically $1$ (see the proof), so that Theorem~\ref{thm:main} gives there
only $e\ge\eF$, while \eqref{eq:localineq} bounds $e-\eF$ from below by an explicit positive definite quadratic form on an explicit
neighbourhood. Near a description of $\Bz$, \eqref{eq:localineq} shows that the deficit itself, and not only the energy excess
$e-\eB$, grows quadratically in the strain and the displacement.

\begin{table}[h]
\begin{tabular}{lllllll}
\toprule
class & $r$ & $\nu$ & $\lambda$ & $B_3$ & $B_4$ & $\varepsilon_0$\\
\midrule
$\Bz$, $\{110\}$ & $83/10000$ & $1/2$ & $0.0077506$ & $0.31143$ & $438.70$ & $3.125\times10^{-5}$\\
                 & $63/10000$ & $2/3$ & $0.0094550$ & $0.52180$ & $827.44$ & $7.186\times10^{-5}$\\
$\Bz$, $\{100\}$ & $23/1000$ & $1/8$ & $0.0094352$ & $0.12806$ & $79.971$ & $9.612\times10^{-6}$\\
                 & $171/10000$ & $1/4$ & $0.0094423$ & $0.15673$ & $166.60$ & $1.140\times10^{-5}$\\
$\Fz$, $\{111\}$ & $179/10000$ & $1/4$ & $0.0088264$ & $0.14248$ & $138.71$ & $1.117\times10^{-5}$\\
                 & $81/5000$ & $1/3$ & $0.0103365$ & $0.17242$ & $212.81$ & $4.808\times10^{-5}$\\
$\Fz$, $\{100\}$ & $11/500$ & $1/6$ & $0.0103365$ & $0.13167$ & $98.235$ & $2.906\times10^{-4}$\\
                 & $187/10000$ & $1/4$ & $0.0103365$ & $0.15174$ & $157.20$ & $4.052\times10^{-5}$\\
\bottomrule
\end{tabular}
\caption{The certified ellipsoids of Theorem~\ref{thm:local}. Here $\lambda$ is a lower bound for the least eigenvalue of the Hessian
at the centre of the function $F$ of \eqref{eq:Fdef} in the coordinates $z=(S,u/\nu)$, $B_3$ and $B_4$ bound its third and fourth
radial derivatives as in \eqref{eq:B3B4}, and $\varepsilon_0$ is the certified lower bound for
$\frac\lambda2-B_3r-\frac1{24}B_4r^2$, computed from the unrounded enclosures. The displayed values of $\lambda$ and $\varepsilon_0$ are
rounded down and those of $B_3$, $B_4$ up.}\label{tab:atlas}
\end{table}

The rest of this section is the proof. Fix a description and a pair $(r,\nu)$, write $s=\|S\|_\Fro$ and $z=(S,u/\nu)\in\R^8$, where
$S$ is expressed in a Frobenius-orthonormal basis, and $e(z):=e(\ltwo Y(S,u))$.

\emph{Criticality.} For every $S$ the configuration $Y(S,0)=A(S)\Lambda_0$ is a lattice, and the point reflection in the origin maps
$Y(S,u)$ to $Y(S,-u)$, because $2w_0\in L_0$. Hence $e(S,u)=e(S,-u)$ and $\partial_ue(S,0)=0$. Let $B_\ast$ be a basis matrix of
$2^{1/3}\Lambda_0$. Then $\ltwo A(S)\Lambda_0=\lone A(S)B_\ast\Z^3$ and $(A(S)B_\ast)^\T A(S)B_\ast=B_\ast^\T C(S)B_\ast$, so that
$e(S,0)=(\Hone(B_\ast^\T C(S)B_\ast)-3)/(2\lone)$ by \eqref{eq:H1}; its differential at $S=0$ vanishes by the criticality shown in
Section~\ref{ssec:strain}. Hence $\nabla e(0)=0$. Moreover $e(0)=\eB$ or $\eF$ by Lemma~\ref{lem:invariance}.

\emph{The defect.} We use the elementary bounds $\|S\|_\op\le\sqrt{2/3}\,s$ and $|\det S|\le s^3/(3\sqrt6)$ for traceless symmetric
$S$, and the identity $\det(I+S)=1-\frac12s^2+\det S$. For $s\le r$ put
\[
  d_-:=1-\tfrac12r^2-\tfrac{r^3}{3\sqrt6},\qquad \mu:=1-\sqrt{\tfrac23}\,r,\qquad \bar\mu:=d_-^{-1/3}\Bigl(1+\sqrt{\tfrac23}\,r\Bigr).
\]
Then $d_-\le\det(I+S)\le1$ and $\mu I\preceq C(S)\preceq\bar\mu I$, so the singular values of $A(S)$ lie in
$[\sqrt\mu,\sqrt{\bar\mu}]$; also $|u|\le r\nu$. For $\Lambda_0=\Fz$ the certificate verifies
$\sqrt{\bar\mu}\,(2^{1/3}+r\nu)\le\frac{51}{40}$, so $D\equiv1$ on $\mathcal E(r,\nu)$ by Lemma~\ref{lem:Dupper}(ii). For
$\Lambda_0=\Bz$ it verifies $\sqrt{\bar\mu}\,(1+r\nu)<\frac{13}{10}$ and $\sqrt{2\mu}-\sqrt{\bar\mu}\,r\nu>\frac{13}{10}$, and
Lemma~\ref{lem:Dupper}(i) with $(\sqrt c-1)^2=(c-1)^2/(1+\sqrt c)^2$ gives
\[
  D\le100\,\|C^{1/2}-I\|_\Fro^2+200\,u^\T Cu\le\frac{100}{(1+\sqrt\mu)^2}\,\|C-I\|_\Fro^2+200\,\bar\mu\,|u|^2 .
\]
Writing $C-I=\vartheta S+(\vartheta-1)I$ with $\vartheta=\det(I+S)^{-1/3}$, we have
$\|C-I\|_\Fro^2=\vartheta^2s^2+3(\vartheta-1)^2$ and, by the mean value theorem,
$0\le\vartheta-1\le\frac13d_-^{-4/3}(\frac12+\frac r{3\sqrt6})s^2$. Hence, on $\mathcal E(r,\nu)$,
\begin{equation}\label{eq:Dquad}
  D\le q_s\,s^2+q_u\,|u|^2,\qquad q_s:=\frac{100\,(d_-^{-2/3}+3b^2)}{(1+\sqrt\mu)^2},\quad q_u:=200\,\bar\mu,\quad
  b:=\tfrac13d_-^{-4/3}\bigl(\tfrac12+\tfrac r{3\sqrt6}\bigr)r .
\end{equation}

\emph{The comparison function.} Let
\begin{equation}\label{eq:Fdef}
  F(z):=\begin{cases}e(z)-\eB-\gap\,\bigl(q_s\,s^2+q_u\,|u|^2\bigr),&\Lambda_0=\Bz,\\ e(z)-\eF,&\Lambda_0=\Fz .\end{cases}
\end{equation}
Then $\Def(\ltwo Y(S,u))\ge F(z)$ on $\mathcal E(r,\nu)$, by \eqref{eq:Dquad} in the first case and by $D\equiv1$ in the second, and
$F(0)=0$, $\nabla F(0)=0$. The certificate provides three numbers $\lambda,B_3,B_4$ with
\begin{equation}\label{eq:B3B4}
  \nabla^2F(0)\succeq\lambda I,\qquad \Bigl|\frac{d^3}{dt^3}F(t\omega)\Big|_{t=0}\Bigr|\le6B_3,\qquad
  \Bigl|\frac{d^4}{dt^4}F(t\omega)\Bigr|\le B_4
\end{equation}
for all unit vectors $\omega\in\R^8$ and all $|t|\le r$. Taylor's formula on the segment from $0$ to $z=t\omega$ gives
\[
  F(z)\ \ge\ \Bigl(\frac\lambda2-B_3r-\frac{B_4r^2}{24}\Bigr)|z|^2 ,
\]
which is \eqref{eq:localineq}. It remains to explain how \eqref{eq:B3B4} is certified; since the subtracted defect majorant is
quadratic, $B_3$ and $B_4$ are bounds for the energy alone.

\emph{The Hessian.} The energy is given by \eqref{eq:H2} with the three families of squared lengths
\[
  (B'm)^\T C\,(B'm),\qquad \bigl(B'(m+f_0)+u\bigr)^\T C\,\bigl(B'(m+f_0)+u\bigr),\qquad (B'^{-\T}m)^\T C^{-1}(B'^{-\T}m),
\]
and the phase $2\pi m\cdot f_0+2\pi(B'^{-\T}m)\cdot u$, where $C^{-1}=\det(I+S)^{-2/3}\adj(I+S)$. The terms with
$\max_i|m_i|\le n_0=8$ are differentiated twice at $z=0$ in ball arithmetic, and the omitted terms are bounded by
Lemma~\ref{lem:tails}. Reflecting the lower triangle gives a symmetric interval matrix $[H]$ containing $\nabla^2F(0)$.
The program first checks that every pivot in the interval $LDL^\T$ factorization of $[H]$ is strictly positive.
It then bisects $[0,2]$ at exact rational midpoints $t$ for $28$ steps. If every interval pivot of $[H]-tI$ is strictly
positive, every enclosed Hessian has least eigenvalue greater than $t$, and $t$ becomes the new lower endpoint.
Otherwise $t$ becomes the upper endpoint of the search. The final lower endpoint is the certified bound $\lambda$.

\emph{The cubic term.} At a two-site description $f_0\in\frac12\Z^3$, so $\cos(2\pi m\cdot f_0)=\pm1$ and the homogeneous cubic part
of the Taylor polynomial at $z=0$ of each term of \eqref{eq:H2} is an explicit polynomial in $(S,u)$ whose coefficients are values of
$\Ga(a+j,\cdot)$, $j\le3$. The certificate sums these polynomials over $\max_i|m_i|\le n_0$ in ball arithmetic. If
$P_3(z)=\sum_{|\beta|=3}c_\beta z^\beta$ is the result, then $|P_3(z)|\le B_{3,\Fro}|z|^3$ with
$B_{3,\Fro}^2=\sum_\beta|c_\beta|^2/\binom3\beta$, the Frobenius norm of the symmetric coefficient tensor, and $B_3$ is $B_{3,\Fro}$ plus
one sixth of the bound of Lemma~\ref{lem:tails} for the third radial derivative of the omitted terms.

\emph{The quartic term.} Along a ray $t\mapsto(tH,tv)$ with $\|H\|_\Fro\le1$, $|v|\le\nu$, we have
$\det(I+tH)=1-\frac12\|H\|_\Fro^2t^2+(\det H)\,t^3$, so the scalar factors $\det(I+tH)^{-1/3}$ and $\det(I+tH)^{-2/3}$ and their first
four $t$-derivatives are enclosed, uniformly in $H$ and in $|t|\le r$, by univariate Taylor arithmetic with the interval parameters
$\|H\|_\Fro^2\in[0,1]$ and $\det H\in[-\frac1{3\sqrt6},\frac1{3\sqrt6}]$. This gives bounds for the operator norms of the first four
derivatives of $C$ and $C^{-1}$ along the ray, hence for the derivatives of each squared length $q$, together with the lower bounds
$q\ge\mu(|B'(m+f_0)|-r\nu)^2$ for the cross-coset terms, where $|B'(m+f_0)|\ge\dist(w_0,L_0)>r\nu$, and their analogues for the other
terms. The fourth derivative of $t\mapsto\Ga(a,\pi q(t))$ is then bounded by Fa\`a di Bruno's formula with absolute values,
\begin{equation}\label{eq:faa4}
  \pi^4\Ga_4q_1^4+6\pi^3\Ga_3q_1^2q_2+\pi^2\Ga_2\,(3q_2^2+4q_1q_3)+\pi\,\Ga_1q_4 ,
\end{equation}
where $q_j\ge|q^{(j)}|$ and $\Ga_j=\Ga(a+j,\pi q_-)$ with $q_-$ the lower bound for $q$; for the reciprocal sum it is combined by
Leibniz's rule with the bounds $|\sin x|\le|x|$, $1-\cos x\le\frac12x^2$ for the phase factor. Summing over $\max_i|m_i|\le n_0$ and
adding the tail bound of Lemma~\ref{lem:tails} gives $B_4$. This estimate uses absolute values termwise and is the reason why the
certified ellipsoids are small.

\begin{remark}\label{rem:optical}
The Hessian at the centre splits into a $5\times5$ strain block, which is the same for all descriptions of a lattice and is governed
by its two shear moduli, and a $3\times3$ block for $u$, which is the dynamical matrix of the lattice at the corresponding
zone-boundary point (Table~\ref{tab:classes}). The positivity of the $3\times3$ blocks at the points $L$ and $X$ of the face-centred
cubic lattice is part of the certificate; it is not a consequence of the stability of the body-centred cubic lattice.
\end{remark}
